\documentclass[11pt]{article}
\usepackage[a4paper,margin=27mm]{geometry}
\usepackage[T1]{fontenc}
\usepackage{lmodern}
\usepackage{amsmath,amssymb}
\usepackage[hidelinks]{hyperref}

\title{Scale Covariance and Pareto Efficiency Do Not\\
  Characterize Nash--Kalai--Smorodinsky Mixtures\\[4pt]
  \large Disproof of Conjecture 00000008249}
\author{Feiyu (Independent researcher)}
\date{26 November 2026}

\begin{document}
\maketitle

\section{The characterization}

Conjecture 00000008249 includes the claim that scale covariance and
Pareto efficiency characterize convex combinations of the Nash and
Kalai--Smorodinsky bargaining solutions. For a two-player bargaining problem
$(S,d)$, write these points as $N(S,d)$ and $K(S,d)$. The characterized rules
would be the maps
\[
  (S,d)\longmapsto \lambda N(S,d)+(1-\lambda)K(S,d),
  \qquad 0\le\lambda\le1 .
\]
We disprove the characterization in both directions. The failure of
sufficiency persists when the weight may depend on $(S,d)$.

On the bargaining domain in \cite[pp.~513--514]{ks}, the set
$S\subset\mathbb R^2$ is compact and convex, every $p\in S$ satisfies $p\ge d$ coordinatewise, and
some $p\in S$ strictly improves both coordinates of $d$. A rule is Pareto
efficient if its value is feasible and no distinct feasible point is
coordinatewise at least as large. Scale covariance means that changing
utility units by $T(x,y)=(a x,b y)$ with $a,b>0$ sends the output for
$(S,d)$ to the output for $(T(S),T(d))$.

\section{A bargaining polygon}

Set $d=(0,0)$ and take
\[
  S=\{(x,y): 0\le x\le32,\quad 0\le y,\quad
                 x+2y\le64,\quad 2x+3y\le108\}.
\]
The set is a closed convex subset of $[0,32]^2$. It contains $d$ and
$(1,1)$, and it is downward comprehensive above $d$ because the upper
bounds have nonnegative coefficients. The points $(32,0)$ and $(0,32)$
show that the coordinate maxima form the ideal point $(32,32)$.

\paragraph{Nash solution.}
Nash maximizes the product of gains over disagreement
\cite[p.~159]{nash}. For every $(x,y)\in S$,
\[
  3xy\le x(108-2x)=1458-2(x-27)^2 .
\]
The feasible point $N=(27,18)$ attains $3xy=1458$. Equality forces
$x=27$, then $y=18$, so the maximum is unique.

\paragraph{Kalai--Smorodinsky solution.}
This solution is the greatest feasible point on the line through
disagreement and the ideal \cite[p.~516]{ks}. Here the line is $y=x$.
Along it, $x+2y\le64$ gives $x\le64/3$, and
$K=(64/3,64/3)$ is feasible. Thus $K$ is the unique line maximum.
The constraints $2x+3y\le108$ and $x+2y\le64$ are tight at $N$ and $K$,
respectively. Since both have positive coefficients, both solutions
are Pareto efficient.

\newpage

\section{The axioms allow a rule outside the segment}

Define $L(S,d)$ by maximizing the first coordinate on $S$, then maximizing
the second coordinate among the points with maximal first coordinate.
Compactness of $S$ gives the first maximum. The first-coordinate maximizing
face is nonempty and closed in $S$, so compactness gives the second maximum.
The two coordinates determine a unique point of $S$.

A feasible point dominating $L(S,d)$ must have the same first coordinate
by the first maximization and the same second coordinate by the second
maximization. Therefore $L$ is Pareto efficient. Each positive affine
change of the two utility coordinates preserves both maximizations, so
$L$ is covariant under such changes. In particular it is scale covariant.

On the bargaining polygon, $L(S,d)=(32,44/3)$ since at $x=32$ the second
upper bound gives $y\le44/3$, and equality is feasible. Every point of
$[K,N]$ has first coordinate at most $27$. Consequently no weight
$\lambda\in[0,1]$ gives $L(S,d)$ on this instance.

\section{Interior mixtures violate Pareto efficiency}

For $0<\lambda<1$ set $(x,y)=\lambda N+(1-\lambda)K$.
This point lies in $S$ by convexity. The two sloping bounds have slacks
\[
  64-(x+2y)=\lambda,\qquad
  108-(2x+3y)=\tfrac43(1-\lambda).
\]
With $\varepsilon=\min\{\lambda,1-\lambda\}/6>0$, increasing both
coordinates by $\varepsilon$ costs $3\varepsilon\le\lambda/2$ in the first
bound and $5\varepsilon\le(5/6)(1-\lambda)$ in the second. The vertical
bound is preserved since $x\le27$ and $\varepsilon\le1/6$. Thus
$(x+\varepsilon,y+\varepsilon)$ is feasible and strictly improves both
utilities. For example, the midpoint $(145/6,59/3)$ is strictly dominated
by $(97/4,79/4)$. Every fixed weight strictly between $0$ and $1$ therefore
fails even the weak Pareto axiom excluding simultaneous strict gains.

\subsection*{Lean formalization}
The project uses Lean and Mathlib \texttt{v4.33.1}. The module
\texttt{Submission.Basic} defines the bargaining domain over
$\mathbb R\times\mathbb R$ with Mathlib's compactness and convexity.
It uses \texttt{IsMaxOn} for Nash and \texttt{IsGreatest} for the
coordinate maxima and the KS line optimum.

\smallskip\noindent
\texttt{disprove\_conjecture\_00000008249} refutes the characterization
of fixed-weight mixtures. The lexicographic witness also refutes the claim
with instance-dependent weights.

\smallskip\noindent
\texttt{interior\_mixture\_strictly\_dominated} verifies the improvement
for every interior weight on the polygon.

\begin{thebibliography}{9}
\bibitem{nash}
J.~F.~Nash, Jr., \emph{The Bargaining Problem},
Econometrica \textbf{18} (1950), 155--162.
\href{https://www.haverford.edu/sites/default/files/Nash1950.pdf}{PDF}.

\bibitem{ks}
E.~Kalai and M.~Smorodinsky, \emph{Other Solutions to Nash's Bargaining
Problem}, Econometrica \textbf{43} (1975), 513--518.
\href{https://www.haverford.edu/sites/default/files/KalaiSmorodinsky1975.pdf}{PDF}.
\end{thebibliography}
\end{document}
